% chapterintro
% ==============================
% 第4章 导言：4.0-chinese.md
本章的主要内容有：

\begin{itemize}
\item 创建一个简单的单元测试

\item 使用Catch2库进行单元测试

\item 使用Google Test库进行单元测试

\item 使用Boost Test进行单元测试

\item 使用动态分析来检测内存缺陷

\item 预期测试失败

\item 使用超时测试运行时间过长的测试

\item 并行测试

\item 运行测试子集

\item 使用测试固件

\end{itemize}

测试代码是开发工具的核心组件。通过单元测试和集成测试自动化测试，不仅可以帮助开发人员尽早回归功能检测，还可以帮助开发人员参与，并了解项目。它可以帮助新开发人员向项目代码提交修改，并确保预期的功能性。对于验证安装是否保留了代码的功能时，自动化测试必不可少。从一开始对单元、模块或库进行测试，可以使用一种纯函数式的风格，将全局变量和全局状态最小化，可让开发者的具有更模块化、更简单的编程风格。

本章中，我们将演示如何使用流行的测试库和框架，将测试集成到CMake构建结构中，并谨记以下目标：

\begin{itemize}
\item 让用户、开发人员和持续集成服务很容易地运行测试集。应该像使用\texttt{Unix Makefile}时，键入\texttt{make test}一样简单。

\item 通过最小化测试时间，高效地运行测试，最大限度地提高运行测试的概率——理想情况下，每次代码修改都该如此。

\end{itemize}


% --- recipe 1: 4.1-chinese.md ---
\section{4.1 创建一个简单的单元测试}

\begin{myInfoBox}
\textbf{NOTE}：此示例代码可以在 \texthref{https://github.com/dev-cafe/cmake-cookbook/tree/v1.0/chapter-04/recipe-01} 中找到，包含一个C++的示例。该示例在CMake 3.5版(或更高版本)中是有效的，并且已经在GNU/Linux、macOS和Windows上进行过测试。
\end{myInfoBox}

CTest是CMake的测试工具，本示例中，我们将使用CTest进行单元测试。为了保持对CMake/CTest的关注，我们的测试代码会尽可能的简单。计划是编写和测试能够对整数求和的代码，示例代码只会对整数进行累加，不处理浮点数。就像年轻的卡尔•弗里德里希•高斯(Carl Friedrich Gauss)，被他的老师测试从1到100求和所有自然数一样，我们将要求代码做同样的事情。为了说明CMake没有对实际测试的语言进行任何限制，我们不仅使用C++可执行文件测试代码，还使用Python脚本和shell脚本作为测试代码。为了简单起见，我们将不使用任何测试库来实现，但是我们将在后面的示例中介绍C++测试框架。

\mySmallsection{准备工作}

代码示例由三个文件组成。实现源文件\texttt{sum\_integs.cpp}对整数向量进行求和，并返回累加结果：

\begin{codebox}[cpp]
#include "sum_integers.hpp"

#include <vector>

int sum_integers(const std::vector<int> integers) {
    auto sum = 0;
    for (auto i : integers) {
        sum += i;
    }
    return sum;
}
\end{codebox}

这个示例是否是优雅的实现并不重要，接口以\texttt{sum\_integers}的形式导出。接口在\texttt{ sum\_integers.hpp }文件中声明，详情如下:

\begin{codebox}[cpp]
#pragma once

#include <vector>

int sum_integers(const std::vector<int> integers);
\end{codebox}

最后，main函数在\texttt{main.cpp}中定义，从\texttt{argv[]}中收集命令行参数，将它们转换成整数向量，调用\texttt{sum\_integers}函数，并将结果打印到输出中:

\begin{codebox}[cpp]
#include "sum_integers.hpp"

#include <iostream>
#include <string>
#include <vector>

// we assume all arguments are integers and we sum them up
// for simplicity we do not verify the type of arguments
int main(int argc, char *argv[]) {
    std::vector<int> integers;
    for (auto i = 1; i < argc; i++) {
        integers.push_back(std::stoi(argv[i]));
    }
    auto sum = sum_integers(integers);
  
    std::cout << sum << std::endl;
}
\end{codebox}

测试这段代码使用C++实现(\texttt{test.cpp})，Bash shell脚本实现(\texttt{test.sh})和Python脚本实现(\texttt{test.py})，只要实现可以返回一个零或非零值，从而CMake可以解释为成功或失败。

C++例子(\texttt{test.cpp})中，我们通过调用\texttt{sum\_integers}来验证1 + 2 + 3 + 4 + 5 = 15：

\begin{codebox}[cpp]
#include "sum_integers.hpp"

#include <vector>

int main() {
    auto integers = {1, 2, 3, 4, 5};
    
  if (sum_integers(integers) == 15) {
        return 0;
    } else {
        return 1;
    }
}
\end{codebox}

Bash shell脚本调用可执行文件：

\begin{codebox}[shell]
#!/usr/bin/env bash

EXECUTABLE=$1

OUTPUT=$($EXECUTABLE 1 2 3 4)

if [ "$OUTPUT" = "10" ]
then
    exit 0
else
    exit 1
fi
\end{codebox}

此外，Python脚本调用可执行文件(使用\texttt{--executable}命令行参数传递)，并使用\texttt{--short}命令行参数执行：

\begin{codebox}[python]
import subprocess
import argparse

# test script expects the executable as argument
parser = argparse.ArgumentParser()
parser.add_argument('--executable',
                                         help='full path to executable')
parser.add_argument('--short',
                                         default=False,
                    action='store_true',
                    help='run a shorter test')
args = parser.parse_args()

def execute_cpp_code(integers):
    result = subprocess.check_output([args.executable] + integers)
    return int(result)

if args.short:
    # we collect [1, 2, ..., 100] as a list of strings
    result = execute_cpp_code([str(i) for i in range(1, 101)])
    assert result == 5050, 'summing up to 100 failed'
else:
    # we collect [1, 2, ..., 1000] as a list of strings
    result = execute_cpp_code([str(i) for i in range(1, 1001)])
    assert result == 500500, 'summing up to 1000 failed'
\end{codebox}

\mySmallsection{具体实施}

现在，我们将逐步描述如何为项目设置测试：

\begin{enumerate}
\item 对于这个例子，我们需要C++11支持，可用的Python解释器，以及Bash shell:

\begin{codebox}[cmake]
cmake_minimum_required(VERSION 3.5 FATAL_ERROR)
project(recipe-01 LANGUAGES CXX)
set(CMAKE_CXX_STANDARD 11)
set(CMAKE_CXX_EXTENSIONS OFF)
set(CMAKE_CXX_STANDARD_REQUIRED ON)
find_package(PythonInterp REQUIRED)
find_program(BASH_EXECUTABLE NAMES bash REQUIRED)
\end{codebox}

\item 然后，定义库及主要可执行文件的依赖关系，以及测试可执行文件：

\begin{codebox}[cmake]
# example library
add_library(sum_integers sum_integers.cpp)
# main code
add_executable(sum_up main.cpp)
target_link_libraries(sum_up sum_integers)
# testing binary
add_executable(cpp_test test.cpp)
target_link_libraries(cpp_test sum_integers)
\end{codebox}

\item 最后，打开测试功能并定义四个测试。最后两个测试， 调用相同的Python脚本，先没有任何命令行参数，再使用\texttt{--short}：

\begin{codebox}[cmake]
enable_testing()
add_test(
  NAME bash_test
  COMMAND ${BASH_EXECUTABLE} ${CMAKE_CURRENT_SOURCE_DIR}/test.sh $<TARGET_FILE:sum_up>
  )
add_test(
  NAME cpp_test
  COMMAND $<TARGET_FILE:cpp_test>
  )
add_test(
  NAME python_test_long
  COMMAND ${PYTHON_EXECUTABLE} ${CMAKE_CURRENT_SOURCE_DIR}/test.py --executable $<TARGET_FILE:sum_up>
  )
add_test(
  NAME python_test_short
  COMMAND ${PYTHON_EXECUTABLE} ${CMAKE_CURRENT_SOURCE_DIR}/test.py --short --executable $<TARGET_FILE:sum_up>
  )
\end{codebox}

\item 现在，我们已经准备好配置和构建代码。先手动进行测试：

\begin{codebox}[shell]
$ mkdir -p build
$ cd build
$ cmake ..
$ cmake --build .
$ ./sum_up 1 2 3 4 5
15
\end{codebox}

\item 然后，我们可以用\texttt{ctest}运行测试集：

\begin{codebox}[shell]
$ ctest
Test project /home/user/cmake-recipes/chapter-04/recipe-01/cxx-example/build
Start 1: bash_test
1/4 Test #1: bash_test ........................ Passed 0.01 sec
Start 2: cpp_test
2/4 Test #2: cpp_test ......................... Passed 0.00 sec
Start 3: python_test_long
3/4 Test #3: python_test_long ................. Passed 0.06 sec
Start 4: python_test_short
4/4 Test #4: python_test_short ................ Passed 0.05 sec
100% tests passed, 0 tests failed out of 4
Total Test time (real) = 0.12 sec
\end{codebox}

\item 还应该尝试中断实现，以验证测试集是否能捕捉到更改。

\end{enumerate}

\mySmallsection{工作原理}

这里的两个关键命令：

\begin{itemize}
\item \texttt{enable\_testing()}，测试这个目录和所有子文件夹(因为我们把它放在主\texttt{CMakeLists.txt})。

\item \texttt{add\_test()}，定义了一个新的测试，并设置测试名称和运行命令。

\end{itemize}

\begin{codebox}[cmake]
add_test(
  NAME cpp_test
  COMMAND $<TARGET_FILE:cpp_test>
  ) 
\end{codebox}

上面的例子中，使用了生成器表达式:\texttt{\$<TARGET\_FILE:cpp\_test>}。生成器表达式，是在生成\textbf{构建系统生成时}的表达式。我们将在第5章第9节中详细地描述生成器表达式。此时，我们可以声明\texttt{ \$<TARGET\_FILE:cpp\_test>}变量，将使用\texttt{cpp\_test}可执行目标的完整路径进行替换。

生成器表达式在测试时非常方便，因为不必显式地将可执行程序的位置和名称，可以硬编码到测试中。以一种可移植的方式实现这一点非常麻烦，因为可执行文件和可执行后缀(例如，Windows上是\texttt{.exe}后缀)的位置在不同的操作系统、构建类型和生成器之间可能有所不同。使用生成器表达式，我们不必显式地了解位置和名称。

也可以将参数传递给要运行的\texttt{test}命令，例如：

\begin{codebox}[cmake]
add_test(
  NAME python_test_short
  COMMAND ${PYTHON_EXECUTABLE} ${CMAKE_CURRENT_SOURCE_DIR}/test.py --short --executable $<TARGET_FILE:sum_up>
  ) 
\end{codebox}

这个例子中，我们按顺序运行测试，并展示如何缩短总测试时间并行执行测试(第8节)，执行测试用例的子集(第9节)。这里，可以自定义测试命令，可以以任何编程语言运行测试集。CTest关心的是，通过命令的返回码测试用例是否通过。CTest遵循的标准约定是，返回零意味着成功，非零返回意味着失败。可以返回零或非零的脚本，都可以做测试用例。

既然知道了如何定义和执行测试，那么了解如何诊断测试失败也很重要。为此，我们可以在代码中引入一个bug，让所有测试都失败:

\begin{codebox}[shell]
Start 1: bash_test
1/4 Test #1: bash_test ........................***Failed 0.01 sec
    Start 2: cpp_test
2/4 Test #2: cpp_test .........................***Failed 0.00 sec
    Start 3: python_test_long
3/4 Test #3: python_test_long .................***Failed 0.06 sec
    Start 4: python_test_short
4/4 Test #4: python_test_short ................***Failed 0.06 sec

0% tests passed, 4 tests failed out of 4

Total Test time (real) = 0.13 sec

The following tests FAILED:
1 - bash_test (Failed)
2 - cpp_test (Failed)
3 - python_test_long (Failed)
4 - python_test_short (Failed)
Errors while running CTest
\end{codebox}

如果我们想了解更多，可以查看文件\texttt{test/Temporary/lasttestsfailure.log}。这个文件包含测试命令的完整输出，并且在分析阶段，要查看的第一个地方。使用以下CLI开关，可以从CTest获得更详细的测试输出：

\begin{itemize}
\item \texttt{--output-on-failure}:将测试程序生成的任何内容打印到屏幕上，以免测试失败。

\item \texttt{-v}:将启用测试的详细输出。

\item \texttt{-vv}:启用更详细的输出。

\end{itemize}

CTest提供了一个非常方快捷的方式，可以重新运行以前失败的测试；要使用的CLI开关是\texttt{--rerun-failed}，在调试期间非常有用。

\mySmallsection{更多信息}

考虑以下定义:

\begin{codebox}[cmake]
add_test(
  NAME python_test_long
  COMMAND ${PYTHON_EXECUTABLE} ${CMAKE_CURRENT_SOURCE_DIR}/test.py --executable $<TARGET_FILE:sum_up>
  )
\end{codebox}

前面的定义可以通过显式指定脚本运行的\texttt{WORKING\_DIRECTORY}重新表达，如下:

\begin{codebox}[cmake]
add_test(
  NAME python_test_long
  COMMAND ${PYTHON_EXECUTABLE} test.py --executable $<TARGET_FILE:sum_up>
  WORKING_DIRECTORY ${CMAKE_CURRENT_SOURCE_DIR}
  )
\end{codebox}

测试名称可以包含\texttt{/}字符，按名称组织相关测试也很有用，例如：

\begin{codebox}[cmake]
add_test(
  NAME python/long
  COMMAND ${PYTHON_EXECUTABLE} test.py --executable $<TARGET_FILE:sum_up>
  WORKING_DIRECTORY ${CMAKE_CURRENT_SOURCE_DIR}
  )
\end{codebox}

有时候，我们需要为测试脚本设置环境变量。这可以通过\texttt{set\_tests\_properties}实现:

\begin{codebox}[cmake]
set_tests_properties(python_test
  PROPERTIES
    ENVIRONMENT
      ACCOUNT_MODULE_PATH=${CMAKE_CURRENT_SOURCE_DIR}
      ACCOUNT_HEADER_FILE=${CMAKE_CURRENT_SOURCE_DIR}/account/account.h
      ACCOUNT_LIBRARY_FILE=$<TARGET_FILE:account>
  )
\end{codebox}

这种方法在不同的平台上并不总可行，CMake提供了解决这个问题的方法。下面的代码片段与上面给出的代码片段相同，在执行实际的Python测试脚本之前，通过\texttt{CMAKE\_COMMAND}调用CMake来预先设置环境变量:

\begin{codebox}[cmake]
add_test(
  NAME
    python_test
  COMMAND
    ${CMAKE_COMMAND} -E env
    ACCOUNT_MODULE_PATH=${CMAKE_CURRENT_SOURCE_DIR}
    ACCOUNT_HEADER_FILE=${CMAKE_CURRENT_SOURCE_DIR}/account/account.h
    ACCOUNT_LIBRARY_FILE=$<TARGET_FILE:account>
    ${PYTHON_EXECUTABLE}
    ${CMAKE_CURRENT_SOURCE_DIR}/account/test.py
  )
\end{codebox}

同样，要注意使用生成器表达式\texttt{ \$<TARGET\_FILE:account> }来传递库文件的位置。

我们已经使用\texttt{ctest}命令执行测试，CMake还将为生成器创建目标(Unix Makefile生成器为\texttt{make test}，Ninja工具为\texttt{ninja test}，或者Visual Studio为\texttt{RUN\_TESTS})。这意味着，还有另一种(几乎)可移植的方法来运行测试：

\begin{codebox}[shell]
$ cmake --build . --target test
\end{codebox}

不幸的是，当使用Visual Studio生成器时，我们需要使用\texttt{RUN\_TESTS}来代替:

\begin{codebox}[shell]
$ cmake --build . --target RUN_TESTS
\end{codebox}

\begin{myInfoBox}
\textbf{NOTE}：\texttt{ctest}提供了丰富的命令行参数。其中一些内容将在以后的示例中探讨。要获得完整的列表，需要使用\texttt{ctest --help}来查看。命令\texttt{cmake --help-manual ctest}会将向屏幕输出完整的ctest手册。
\end{myInfoBox}


% --- recipe 2: 4.2-chinese.md ---
\section{4.2 使用Catch2库进行单元测试}

\begin{myInfoBox}
\textbf{NOTE}：此示例代码可以在 \texthref{https://github.com/dev-cafe/cmake-cookbook/tree/v1.0/chapter-04/recipe-02} 中找到，包含一个C++的示例。该示例在CMake 3.5版(或更高版本)中是有效的，并且已经在GNU/Linux、macOS和Windows上进行过测试。
\end{myInfoBox}

前面的配置中，使用返回码来表示\texttt{test.cpp}测试的成功或失败。对于简单功能没问题，但是通常情况下，我们想要使用一个测试框架，它提供了相关基础设施来运行更复杂的测试，包括固定方式进行测试，与数值公差的比较，以及在测试失败时输出更好的错误报告。这里，我们用目前比较流行的测试库Catch2( \texthref{https://github.com/catchorg/Catch2} )来进行演示。这个测试框架有个很好的特性，它可以通过单个头库包含在项目中进行测试，这使得编译和更新框架特别容易。这个配置中，我们将CMake和Catch2结合使用，来测试上一个求和代码。

我们需要\texttt{catch.hpp}头文件，可以从 \texthref{https://github.com/catchorg/Catch2} (我们使用的是版本2.0.1)下载，并将它与\texttt{test.cpp}一起放在项目的根目录下。

\mySmallsection{准备工作}

\texttt{main.cpp}、\texttt{sum\_integers.cpp}和\texttt{sum\_integers.hpp}与之前的示例相同，但将更新\texttt{test.cpp}:

\begin{codebox}[cpp]
#include "sum_integers.hpp"

// this tells catch to provide a main()
// only do this in one cpp file
#define CATCH_CONFIG_MAIN
#include "catch.hpp"
#include <vector>

TEST_CASE("Sum of integers for a short vector", "[short]")
{
  auto integers = {1, 2, 3, 4, 5};
  REQUIRE(sum_integers(integers) == 15);
}

TEST_CASE("Sum of integers for a longer vector", "[long]")
{
  std::vector<int> integers;
  for (int i = 1; i < 1001; ++i)
  {
    integers.push_back(i);
  }
  REQUIRE(sum_integers(integers) == 500500);
}
\end{codebox}

\texttt{catch.hpp}头文件可以从\texthref{https://github.com/catchorg/Catch2} (版本为2.0.1)下载，并将它与\texttt{test.cpp}放在项目的根目录中。

\mySmallsection{具体实施}

使用Catch2库，需要修改之前的所使用\texttt{CMakeList.txt}：

\begin{enumerate}
\item 保持\texttt{CMakeLists.txt}大多数部分内容不变:

\begin{codebox}[cmake]
# set minimum cmake version
cmake_minimum_required(VERSION 3.5 FATAL_ERROR)
# project name and language
project(recipe-02 LANGUAGES CXX)
# require C++11
set(CMAKE_CXX_STANDARD 11)
set(CMAKE_CXX_EXTENSIONS OFF)
set(CMAKE_CXX_STANDARD_REQUIRED ON)
# example library
add_library(sum_integers sum_integers.cpp)
# main code
add_executable(sum_up main.cpp)
target_link_libraries(sum_up sum_integers)
# testing binary
add_executable(cpp_test test.cpp)
target_link_libraries(cpp_test sum_integers)
\end{codebox}

\item 对于上一个示例的配置，需要保留一个测试，并重命名它。注意，\texttt{--success}选项可传递给单元测试的可执行文件。这是一个Catch2选项，测试成功时，也会有输出:

\begin{codebox}[cmake]
enable_testing()
add_test(
  NAME catch_test
  COMMAND $<TARGET_FILE:cpp_test> --success
  )
\end{codebox}

\item 就是这样！让我们来配置、构建和测试。CTest中，使用\texttt{-V}选项运行测试，以获得单元测试可执行文件的输出:

\begin{codebox}[shell]
$ mkdir -p build
$ cd build
$ cmake ..
$ cmake --build .
$ ctest -V
UpdateCTestConfiguration from :/home/user/cmake-cookbook/chapter-04/recipe-02/cxx-example/build/DartConfiguration.tcl
UpdateCTestConfiguration from :/home/user/cmake-cookbook/chapter-04/recipe-02/cxx-example/build/DartConfiguration.tcl
Test project /home/user/cmake-cookbook/chapter-04/recipe-02/cxx-example/build
Constructing a list of tests
Done constructing a list of tests
Updating test list for fixtures
Added 0 tests to meet fixture requirements
Checking test dependency graph...
Checking test dependency graph end
test 1
Start 1: catch_test
1: Test command: /home/user/cmake-cookbook/chapter-04/recipe-02/cxx-example/build/cpp_test "--success"
1: Test timeout computed to be: 10000000
1:
1: ~~~~~~~~~~~~~~~~~~~~~~~~~~~~~~~~~~~~~~~~~~~~~~~~~~~~~~~~~~~~~~~~~~~
1: cpp_test is a Catch v2.0.1 host application.
1: Run with -? for options
1:
1: ----------------------------------------------------------------
1: Sum of integers for a short vector
1: ----------------------------------------------------------------
1: /home/user/cmake-cookbook/chapter-04/recipe-02/cxx-example/test.cpp:10
1: ...................................................................
1:
1: /home/user/cmake-cookbook/chapter-04/recipe-02/cxx-example/test.cpp:12:
1: PASSED:
1: REQUIRE( sum_integers(integers) == 15 )
1: with expansion:
1: 15 == 15
1:
1: ----------------------------------------------------------------
1: Sum of integers for a longer vector
1: ----------------------------------------------------------------
1: /home/user/cmake-cookbook/chapter-04/recipe-02/cxx-example/test.cpp:15
1: ...................................................................
1:
1: /home/user/cmake-cookbook/chapter-04/recipe-02/cxx-example/test.cpp:20:
1: PASSED:
1: REQUIRE( sum_integers(integers) == 500500 )
1: with expansion:
1: 500500 (0x7a314) == 500500 (0x7a314)
1:
1: ===================================================================
1: All tests passed (2 assertions in 2 test cases)
1:
1/1 Test #1: catch_test ....................... Passed 0.00 s
100% tests passed, 0 tests failed out of 1
Total Test time (real) = 0.00 se
\end{codebox}

\item 我们也可以测试\texttt{cpp\_test}的二进制文件，可以直接从Catch2中看到输出:

\begin{codebox}[shell]
$ ./cpp_test --success
~~~~~~~~~~~~~~~~~~~~~~~~~~~~~~~~~~~~~~~~~~~~~~~~~~~~~~~~~~~~~~~~~~~
cpp_test is a Catch v2.0.1 host application.
Run with -? for options
-------------------------------------------------------------------
Sum of integers for a short vector
-------------------------------------------------------------------
/home/user/cmake-cookbook/chapter-04/recipe-02/cxx-example/test.cpp:10
...................................................................
/home/user/cmake-cookbook/chapter-04/recipe-02/cxx-example/test.cpp:12:
PASSED:
REQUIRE( sum_integers(integers) == 15 )
with expansion:
15 == 15
-------------------------------------------------------------------
Sum of integers for a longer vector
-------------------------------------------------------------------
/home/user/cmake-cookbook/chapter-04/recipe-02/cxx-example/test.cpp:15
...................................................................
/home/user/cmake-cookbook/chapter-04/recipe-02/cxx-example/test.cpp:20:
PASSED:
REQUIRE( sum_integers(integers) == 500500 )
with expansion:
500500 (0x7a314) == 500500 (0x7a314)
===================================================================
All tests passed (2 assertions in 2 test cases)
\end{codebox}

\item Catch2将生成一个可执行文件，还可以尝试执行以下命令，以探索单元测试框架提供的选项:

\begin{codebox}[shell]
$ ./cpp_test --help
\end{codebox}

\end{enumerate}

\mySmallsection{工作原理}

Catch2是一个单头文件测试框架，所以不需要定义和构建额外的目标。只需要确保CMake能找到\texttt{catch.hpp}，从而构建\texttt{test.cpp}即可。为了方便起见，将它放在与\texttt{test.cpp}相同的目录中，我们可以选择一个不同的位置，并使用\texttt{target\_include\_directory}指示该位置。另一种方法是将头部封装到接口库中，这可以在Catch2文档中说明( \texthref{https://github.com/catchorg/catch2/blob/maste/docs/build.systems.md\#cmake} ):

\begin{codebox}[cmake]
# Prepare "Catch" library for other executables 
set(CATCH_INCLUDE_DIR
${CMAKE_CURRENT_SOURCE_DIR}/catch) 

add_library(Catch
INTERFACE) 

target_include_directories(Catch INTERFACE
${CATCH_INCLUDE_DIR})
\end{codebox}

然后，我们对库进行如下链接:

\begin{codebox}[cmake]
target_link_libraries(cpp_test Catch)
\end{codebox}

回想一下第3中的讨论，在第1章从简单的可执行库到接口库，是CMake提供的伪目标库，这些伪目标库对于指定项目外部目标的需求非常有用。

\mySmallsection{更多信息}

这是一个简单的例子，主要关注CMake。当然，Catch2提供了更多功能。有关Catch2框架的完整文档，可访问 \texthref{https://github.com/catchorg/Catch2} 。

Catch2代码库包含有CMake函数，用于解析Catch测试并自动创建CMake测试，不需要显式地输入\texttt{add\_test()}函数，可见 \texthref{https://github.com/catchorg/Catch2/blob/master/contrib/ParseAndAddCatchTests.cmake} 。


% --- recipe 3: 4.3-chinese.md ---
\section{4.3 使用Google Test库进行单元测试}

\begin{myInfoBox}
\textbf{NOTE}：此示例代码可以在 \texthref{https://github.com/dev-cafe/cmake-cookbook/tree/v1.0/chapter-04/recipe-03} 中找到，包含一个C++的示例。该示例在CMake 3.11版(或更高版本)中是有效的，并且已经在GNU/Linux、macOS和Windows上进行过测试。在代码库中，有一个支持CMake 3.5的例子。
\end{myInfoBox}

本示例中，我们将演示如何在CMake的帮助下使用Google Test框架实现单元测试。与前一个配置相比，Google Test框架不仅仅是一个头文件，也是一个库，包含两个需要构建和链接的文件。可以将它们与我们的代码项目放在一起，但是为了使代码项目更加轻量级，我们将选择在配置时，下载一个定义良好的Google Test，然后构建框架并链接它。我们将使用较新的\texttt{FetchContent}模块(从CMake版本3.11开始可用)。第8章中会继续讨论\texttt{FetchContent}，在这里将讨论模块在底层是如何工作的，并且还将演示如何使用\texttt{ExternalProject\_Add}进行模拟。此示例的灵感来自(改编自) \texthref{https://cmake.org/cmake/help/v3.11/module/FetchContent.html} 示例。

\mySmallsection{准备工作}

\texttt{main.cpp}、\texttt{sum\_integers.cpp}和\texttt{sum\_integers.hpp}与之前相同，修改\texttt{test.cpp}:

\begin{codebox}[cpp]
#include "sum_integers.hpp"
#include "gtest/gtest.h"

#include <vector>

int main(int argc, char **argv) {
  ::testing::InitGoogleTest(&argc, argv);
  return RUN_ALL_TESTS();
}

TEST(example, sum_zero) {
  auto integers = {1, -1, 2, -2, 3, -3};
  auto result = sum_integers(integers);
  ASSERT_EQ(result, 0);
}

TEST(example, sum_five) {
  auto integers = {1, 2, 3, 4, 5};
  auto result = sum_integers(integers);
  ASSERT_EQ(result, 15);
}
\end{codebox}

如上面的代码所示，我们显式地将\texttt{gtest.h}，而不将其他Google Test源放在代码项目存储库中，会在配置时使用\texttt{FetchContent}模块下载它们。

\mySmallsection{具体实施}

下面的步骤描述了如何设置\texttt{CMakeLists.txt}，使用GTest编译可执行文件及其相应的测试:

\begin{enumerate}
\item 与前两个示例相比，\texttt{CMakeLists.txt}的开头基本没有变化，CMake 3.11才能使用\texttt{FetchContent}模块:

\begin{codebox}[cmake]
# set minimum cmake version
cmake_minimum_required(VERSION 3.11 FATAL_ERROR)
# project name and language
project(recipe-03 LANGUAGES CXX)
# require C++11
set(CMAKE_CXX_STANDARD 11)
set(CMAKE_CXX_EXTENSIONS OFF)
set(CMAKE_CXX_STANDARD_REQUIRED ON)
set(CMAKE_WINDOWS_EXPORT_ALL_SYMBOLS ON)
# example library
add_library(sum_integers sum_integers.cpp)
# main code
add_executable(sum_up main.cpp)
target_link_libraries(sum_up sum_integers)
\end{codebox}

\item 然后引入一个\texttt{if}，检查\texttt{ENABLE\_UNIT\_TESTS}。默认情况下，它为\texttt{ON}，但有时需要设置为\texttt{OFF}，以免在没有网络连接时，也能使用Google Test:

\begin{codebox}[cmake]
option(ENABLE_UNIT_TESTS "Enable unit tests" ON)
message(STATUS "Enable testing: ${ENABLE_UNIT_TESTS}")
if(ENABLE_UNIT_TESTS)
    # all the remaining CMake code will be placed here
endif()
\end{codebox}

\item \texttt{if}内部包含\texttt{FetchContent}模块，声明要获取的新内容，并查询其属性:

\begin{codebox}[cmake]
include(FetchContent)
FetchContent_Declare(
  googletest
  GIT_REPOSITORY https://github.com/google/googletest.git
  GIT_TAG release-1.8.0
)
FetchContent_GetProperties(googletest)
\end{codebox}

\item 如果内容还没有获取到，将尝试获取并配置它。这需要添加几个可以链接的目标。本例中，我们对\texttt{gtest\_main}感兴趣。该示例还包含一些变通方法，用于使用在Visual Studio下的编译:

\begin{codebox}[cmake]
if(NOT googletest_POPULATED)
  FetchContent_Populate(googletest)
  # Prevent GoogleTest from overriding our compiler/linker options
  # when building with Visual Studio
  set(gtest_force_shared_crt ON CACHE BOOL "" FORCE)
  # Prevent GoogleTest from using PThreads
  set(gtest_disable_pthreads ON CACHE BOOL "" FORCE)
  # adds the targers: gtest, gtest_main, gmock, gmock_main
  add_subdirectory(
    ${googletest_SOURCE_DIR}
    ${googletest_BINARY_DIR}
    )
  # Silence std::tr1 warning on MSVC
  if(MSVC)
    foreach(_tgt gtest gtest_main gmock gmock_main)
      target_compile_definitions(${_tgt}
        PRIVATE
            "_SILENCE_TR1_NAMESPACE_DEPRECATION_WARNING"
      )
    endforeach()
  endif()
endif()
\end{codebox}

\item 然后，使用\texttt{target\_sources}和\texttt{target\_link\_libraries}命令，定义\texttt{cpp\_test}可执行目标并指定它的源文件:

\begin{codebox}[cmake]
add_executable(cpp_test "")
target_sources(cpp_test
  PRIVATE
    test.cpp
  )
target_link_libraries(cpp_test
  PRIVATE
    sum_integers
    gtest_main
  )
\end{codebox}

\item 最后，使用\texttt{enable\_test}和\texttt{add\_test}命令来定义单元测试:

\begin{codebox}[cmake]
enable_testing()
add_test(
  NAME google_test
  COMMAND $<TARGET_FILE:cpp_test>
  )
\end{codebox}

\item 现在，准备配置、构建和测试项目:

\begin{codebox}[shell]
$ mkdir -p build
$ cd build
$ cmake ..
$ cmake --build .
$ ctest
Test project /home/user/cmake-cookbook/chapter-04/recipe-03/cxx-example/build
    Start 1: google_test
1/1 Test #1: google_test ...................... Passed 0.00 sec
100% tests passed, 0 tests failed out of 1
Total Test time (real) = 0.00 sec
\end{codebox}

\item 可以直接运行\texttt{cpp\_test}:

\begin{codebox}[shell]
$ ./cpp_test
[==========] Running 2 tests from 1 test case.
[----------] Global test environment set-up.
[----------] 2 tests from example
[ RUN ] example.sum_zero
[ OK ] example.sum_zero (0 ms)
[ RUN ] example.sum_five
[ OK ] example.sum_five (0 ms)
[----------] 2 tests from example (0 ms total)
[----------] Global test environment tear-down
[==========] 2 tests from 1 test case ran. (0 ms total)
[ PASSED ] 2 tests.
\end{codebox}

\end{enumerate}

\mySmallsection{工作原理}

\texttt{FetchContent}模块支持通过\texttt{ExternalProject}模块，在配置时填充内容，并在其3.11版本中成为CMake的标准部分。而\texttt{ExternalProject\_Add()}在构建时(见第8章)进行下载操作，这样\texttt{FetchContent}模块使得构建可以立即进行，这样获取的主要项目和外部项目(在本例中为Google Test)仅在第一次执行CMake时调用，使用\texttt{add\_subdirectory}可以嵌套。

为了获取Google Test，首先声明外部内容:

\begin{codebox}[cmake]
include(FetchContent)

FetchContent_Declare(
    googletest
  GIT_REPOSITORY https://github.com/google/googletest.git
  GIT_TAG release-1.8.0
)
\end{codebox}

本例中，我们获取了一个带有特定标记的Git库(release-1.8.0)，但是我们也可以从Subversion、Mercurial或HTTP(S)源获取一个外部项目。有关可用选项，可参考相应的\texttt{ExternalProject\_Add}命令的选项，网址是\texthref{https://cmake.org/cmake/help/v3.11/module/ExternalProject.html} 。

调用\texttt{FetchContent\_Populate()}之前，检查是否已经使用\texttt{FetchContent\_GetProperties()}命令处理了内容填充；否则，调用\texttt{FetchContent\_Populate()}超过一次后，就会抛出错误。

\texttt{FetchContent\_Populate(googletest)}用于填充源并定义\texttt{googletest\_SOURCE\_DIR}和\texttt{googletest\_BINARY\_DIR}，可以使用它们来处理Google Test项目(使用\texttt{add\_subdirectory()}，因为它恰好也是一个CMake项目):

\begin{codebox}[cmake]
add_subdirectory(
  ${googletest_SOURCE_DIR}
  ${googletest_BINARY_DIR}
  )
\end{codebox}

前面定义了以下目标：\texttt{gtest}、\texttt{gtest\_main}、\texttt{gmock}和\texttt{gmock\_main}。这个配置中，作为单元测试示例的库依赖项，我们只对\texttt{gtest\_main}目标感兴趣：

\begin{codebox}[cmake]
target_link_libraries(cpp_test
  PRIVATE
    sum_integers
    gtest_main
)
\end{codebox}

构建代码时，可以看到如何正确地对Google Test进行配置和构建。有时，我们希望升级到更新的Google Test版本，这时需要更改的唯一一行就是详细说明\texttt{GIT\_TAG}的那一行。

\mySmallsection{更多信息}

了解了\texttt{FetchContent}及其构建时的近亲\texttt{ExternalProject\_Add}，我们将在第8章中重新讨论这些命令。有关可用选项的详细讨论，可参考\texthref{https://cmake.org/cmake/help/v3.11/module/FetchContent.html} 。

本示例中，我们在配置时获取源代码，也可以将它们安装在系统环境中，并使用\texttt{FindGTest}模块来检测库和头文件(https://cmake.org/cmake/help/v3.5/module/FindTest.html )。从3.9版开始，CMake还提供了一个Google Test模块(https://cmake.org/cmake/help/v3.9/module/GoogleTest.html )，它提供了一个\texttt{gtest\_add\_tests}函数。通过搜索Google Test宏的源代码，可以使用此函数自动添加测试。

当然，Google Test有许多有趣的的特性，可在 \texthref{https://github.com/google/googletest} 查看。


% --- recipe 4: 4.4-chinese.md ---
\section{4.4 使用Boost Test进行单元测试}

\begin{myInfoBox}
\textbf{NOTE}：此示例代码可以在 \texthref{https://github.com/dev-cafe/cmake-cookbook/tree/v1.0/chapter-04/recipe-04} 中找到，包含一个C++的示例。该示例在CMake 3.5版(或更高版本)中是有效的，并且已经在GNU/Linux、macOS和Windows上进行过测试。
\end{myInfoBox}

Boost Test是在C++社区中，一个非常流行的单元测试框架。本例中，我们将演示如何使用Boost Test，对求和示例代码进行单元测试。

\mySmallsection{准备工作}

\texttt{main.cpp}、\texttt{sum\_integers.cpp}和\texttt{sum\_integers.hpp}与之前的示例相同，将更新\texttt{test.cpp}作为使用Boost Test库进行的单元测试：

\begin{codebox}[cpp]
#include "sum_integers.hpp"

#include <vector>

#define BOOST_TEST_MODULE example_test_suite
#include <boost/test/unit_test.hpp>
BOOST_AUTO_TEST_CASE(add_example)
{
  auto integers = {1, 2, 3, 4, 5};
  auto result = sum_integers(integers);
  BOOST_REQUIRE(result == 15);
}
\end{codebox}

\mySmallsection{具体实施}

以下是使用Boost Test构建项目的步骤:

\begin{enumerate}
\item 先从\texttt{CMakeLists.txt}开始:

\begin{codebox}[cmake]
# set minimum cmake version
cmake_minimum_required(VERSION 3.5 FATAL_ERROR)
# project name and language
project(recipe-04 LANGUAGES CXX)
# require C++11
set(CMAKE_CXX_STANDARD 11)
set(CMAKE_CXX_EXTENSIONS OFF)
set(CMAKE_CXX_STANDARD_REQUIRED ON)
# example library
add_library(sum_integers sum_integers.cpp)
# main code
add_executable(sum_up main.cpp)
target_link_libraries(sum_up sum_integers)
\end{codebox}

\item 检测Boost库并将\texttt{cpp\_test}链接到它:

\begin{codebox}[cmake]
find_package(Boost 1.54 REQUIRED COMPONENTS unit_test_framework)
add_executable(cpp_test test.cpp)
target_link_libraries(cpp_test
  PRIVATE
    sum_integers
    Boost::unit_test_framework
  )
# avoid undefined reference to "main" in test.cpp
target_compile_definitions(cpp_test
  PRIVATE
    BOOST_TEST_DYN_LINK
  )
\end{codebox}

\item 最后，定义单元测试:

\begin{codebox}[cmake]
enable_testing()
add_test(
  NAME boost_test
  COMMAND $<TARGET_FILE:cpp_test>
  )
\end{codebox}

\item 下面是需要配置、构建和测试代码的所有内容:

\begin{codebox}[shell]
$ mkdir -p build
$ cd build
$ cmake ..
$ cmake --build .
$ ctest
Test project /home/user/cmake-recipes/chapter-04/recipe-04/cxx-example/build
Start 1: boost_test
1/1 Test #1: boost_test ....................... Passed 0.01 sec
100% tests passed, 0 tests failed out of 1
Total Test time (real) = 0.01 sec
$ ./cpp_test
Running 1 test case...
*** No errors detected
\end{codebox}

\end{enumerate}

\mySmallsection{工作原理}

使用\texttt{find\_package}来检测Boost的\texttt{unit\_test\_framework}组件(参见第3章，第8节)。我们认为这个组件是\texttt{REQUIRED}的，如果在系统环境中找不到它，配置将停止。\texttt{cpp\_test}目标需要知道在哪里可以找到Boost头文件，并且需要链接到相应的库；它们都由\texttt{IMPORTED}库目标\texttt{Boost::unit\_test\_framework}提供，该目标由\texttt{find\_package}设置。

\mySmallsection{更多信息}

本示例中，我们假设系统上安装了Boost。或者，我们可以在编译时获取并构建Boost依赖项。然而，Boost不是轻量级依赖项。我们的示例代码中，我们只使用了最基本的设施，但是Boost提供了丰富的特性和选项，有感兴趣的读者可以去这里看看：\texthref{http://www.boost.org/doc/libs/1\_65\_1/libs/test/doc/html/index.html} 。


% --- recipe 5: 4.5-chinese.md ---
\section{4.5 使用动态分析来检测内存缺陷}

\begin{myInfoBox}
\textbf{NOTE}：此示例代码可以在 \texthref{https://github.com/dev-cafe/cmake-cookbook/tree/v1.0/chapter-04/recipe-05} 中找到，包含一个C++的示例。该示例在CMake 3.5版(或更高版本)中是有效的，并且已经在GNU/Linux、macOS和Windows上进行过测试。
\end{myInfoBox}

内存缺陷：写入或读取越界，或者内存泄漏(已分配但从未释放的内存)，会产生难以跟踪的bug，最好尽早将它们检查出来。Valgrind( \texthref{http://valgrind.org} )是一个通用的工具，用来检测内存缺陷和内存泄漏。本节中，我们将在使用CMake/CTest测试时使用Valgrind对内存问题进行警告。

\mySmallsection{准备工作}

对于这个配置，需要三个文件。第一个是测试的实现(我们可以调用文件\texttt{leaky\_implementation.cpp}):

\begin{codebox}[cpp]
#include "leaky_implementation.hpp"

int do_some_work() {
  
  // we allocate an array
  double *my_array = new double[1000];
  
  // do some work
  // ...
  
  // we forget to deallocate it
  // delete[] my_array;
  
  return 0;
}
\end{codebox}

还需要相应的头文件(\texttt{leaky\_implementation.hpp}):

\begin{codebox}[cpp]
#pragma once

int do_some_work();
\end{codebox}

并且，需要测试文件(\texttt{test.cpp}):

\begin{codebox}[cpp]
#include "leaky_implementation.hpp"

int main() {
  int return_code = do_some_work();
  
  return return_code;
}
\end{codebox}

我们希望测试通过，因为\texttt{return\_code}硬编码为\texttt{0}。这里我们也期望检测到内存泄漏，因为\texttt{my\_array}没有释放。

\mySmallsection{具体实施}

下面展示了如何设置CMakeLists.txt来执行代码动态分析:

\begin{enumerate}
\item 我们首先定义CMake最低版本、项目名称、语言、目标和依赖关系:

\begin{codebox}[cmake]
cmake_minimum_required(VERSION 3.5 FATAL_ERROR)
project(recipe-05 LANGUAGES CXX)
set(CMAKE_CXX_STANDARD 11)
set(CMAKE_CXX_EXTENSIONS OFF)
set(CMAKE_CXX_STANDARD_REQUIRED ON)
add_library(example_library leaky_implementation.cpp)
add_executable(cpp_test test.cpp)
target_link_libraries(cpp_test example_library)
\end{codebox}

\item 然后，定义测试目标，还定义了\texttt{MEMORYCHECK\_COMMAND}:

\begin{codebox}[cmake]
find_program(MEMORYCHECK_COMMAND NAMES valgrind)
set(MEMORYCHECK_COMMAND_OPTIONS "--trace-children=yes --leak-check=full")
# add memcheck test action
include(CTest)
enable_testing()
add_test(
  NAME cpp_test
  COMMAND $<TARGET_FILE:cpp_test>
  )
\end{codebox}

\item 运行测试集，报告测试通过情况，如下所示:

\begin{codebox}[shell]
$ ctest
Test project /home/user/cmake-recipes/chapter-04/recipe-05/cxx-example/build
Start 1: cpp_test
1/1 Test #1: cpp_test ......................... Passed 0.00 sec
100% tests passed, 0 tests failed out of 1
Total Test time (real) = 0.00 sec
\end{codebox}

\item 现在，我们希望检查内存缺陷，可以观察到被检测到的内存泄漏:

\begin{codebox}[shell]
$ ctest -T memcheck
Site: myhost
Build name: Linux-c++
Create new tag: 20171127-1717 - Experimental
Memory check project /home/user/cmake-recipes/chapter-04/recipe-05/cxx-example/build
Start 1: cpp_test
1/1 MemCheck #1: cpp_test ......................... Passed 0.40 sec
100% tests passed, 0 tests failed out of 1
Total Test time (real) = 0.40 sec
-- Processing memory checking output:
1/1 MemCheck: #1: cpp_test ......................... Defects: 1
MemCheck log files can be found here: ( * corresponds to test number)
/home/user/cmake-recipes/chapter-04/recipe-05/cxx-example/build/Testing/Temporary/MemoryChecker.*.log
Memory checking results:
Memory Leak - 1
\end{codebox}

\item 最后一步，应该尝试修复内存泄漏，并验证\texttt{ctest -T memcheck}没有报告错误。

\end{enumerate}

\mySmallsection{工作原理}

使用\texttt{find\_program(MEMORYCHECK\_COMMAND NAMES valgrind)}查找valgrind，并将\texttt{MEMORYCHECK\_COMMAND}设置为其绝对路径。我们显式地包含CTest模块来启用\texttt{memcheck}测试操作，可以使用\texttt{CTest -T memcheck}来启用这个操作。此外，使用\texttt{set(MEMORYCHECK\_COMMAND\_OPTIONS "--trace-children=yes --leak-check=full")}，将相关参数传递给Valgrind。内存检查会创建一个日志文件，该文件可用于详细记录内存缺陷信息。

\begin{myInfoBox}
\textbf{NOTE}：一些工具，如代码覆盖率和静态分析工具，可以进行类似地设置。然而，其中一些工具的使用更加复杂，因为需要专门的构建和工具链。Sanitizers就是这样一个例子。有关更多信息，请参见\texthref{https://github.com/arsenm/sanitizers-cmake} 。另外，请参阅第14章，其中讨论了AddressSanitizer和ThreadSanitizer。
\end{myInfoBox}

\mySmallsection{更多信息}

该方法可向测试面板报告内存缺陷，这里演示的功能也可以独立于测试面板使用。我们将在第14章中重新讨论，与CDash一起使用的情况。

有关Valgrind及其特性和选项的文档，请参见\texthref{http://valgrind.org} 。


% --- recipe 6: 4.6-chinese.md ---
\section{4.6 预期测试失败}

\begin{myInfoBox}
\textbf{NOTE}：此示例代码可以在 \texthref{https://github.com/dev-cafe/cmake-cookbook/tree/v1.0/chapter-04/recipe-06} 中找到，包含一个C++的示例。该示例在CMake 3.5版(或更高版本)中是有效的，并且已经在GNU/Linux、macOS和Windows上进行过测试。
\end{myInfoBox}

理想情况下，我们希望所有的测试能在每个平台上通过。然而，也可能想要测试预期的失败或异常是否会在受控的设置中进行。这种情况下，我们将把预期的失败定义为成功。我们认为，这通常应该交给测试框架(例如：Catch2或Google Test)的任务，它应该检查预期的失败并向CMake报告成功。但是，在某些情况下，您可能希望将测试的非零返回代码定义为成功；换句话说，您可能想要颠倒成功和失败的定义。在本示例中，我们将演示这种情况。

\mySmallsection{准备工作}

这个配置的测试用例是一个很小的Python脚本(\texttt{test.py})，它总是返回1，CMake将其解释为失败:

\begin{codebox}[python]
import sys

# simulate a failing test
sys.exit(1)
\end{codebox}

\mySmallsection{实施步骤}

如何编写CMakeLists.txt来完成我们的任务:

\begin{enumerate}
\item 这个示例中，不需要任何语言支持从CMake，但需要Python:

\begin{codebox}[cmake]
cmake_minimum_required(VERSION 3.5 FATAL_ERROR)
project(recipe-06 LANGUAGES NONE)
find_package(PythonInterp REQUIRED)
\end{codebox}

\item 然后，定义测试并告诉CMake，测试预期会失败:

\begin{codebox}[cmake]
enable_testing()
add_test(example ${PYTHON_EXECUTABLE} ${CMAKE_CURRENT_SOURCE_DIR}/test.py)
set_tests_properties(example PROPERTIES WILL_FAIL true)
\end{codebox}

\item 最后，报告是一个成功的测试，如下所示:

\begin{codebox}[shell]
$ mkdir -p build
$ cd build
$ cmake ..
$ cmake --build .
$ ctest
Test project /home/user/cmake-recipes/chapter-04/recipe-06/example/build
Start 1: example
1/1 Test #1: example .......................... Passed 0.00 sec
100% tests passed, 0 tests failed out of 1
Total Test time (real) = 0.01 sec
\end{codebox}

\end{enumerate}

\mySmallsection{工作原理}

使用\texttt{set\_tests\_properties(example PROPERTIES WILL\_FAIL true)}，将属性\texttt{WILL\_FAIL}设置为\texttt{true}，这将转换成功与失败。但是，这个特性不应该用来临时修复损坏的测试。

\mySmallsection{更多信息}

如果需要更大的灵活性，可以将测试属性\texttt{PASS\_REGULAR\_EXPRESSION}和\texttt{FAIL\_REGULAR\_EXPRESSION}与\texttt{set\_tests\_properties}组合使用。如果设置了这些参数，测试输出将根据参数给出的正则表达式列表进行检查，如果匹配了正则表达式，测试将通过或失败。可以在测试中设置其他属性，完整的属性列表可以参考：\texthref{https://cmake.org/cmake/help/v3.5/manual/cmake-properties.7.html\#properties-on-tests} 。


% --- recipe 7: 4.7-chinese.md ---
\section{4.7 使用超时测试运行时间过长的测试}

\begin{myInfoBox}
\textbf{NOTE}：此示例代码可以在 \texthref{https://github.com/dev-cafe/cmake-cookbook/tree/v1.0/chapter-04/recipe-07} 中找到。该示例在CMake 3.5版(或更高版本)中是有效的，并且已经在GNU/Linux、macOS和Windows上进行过测试。
\end{myInfoBox}

理想情况下，测试集应该花很短的时间进行，以便开发人员经常运行测试，并使每个提交(变更集)进行测试成为可能(或更容易)。然而，有些测试可能会花费更长的时间或者被卡住(例如，由于高文件I/O负载)，我们可能需要设置超时来终止耗时过长的测试，它们延迟了整个测试，并阻塞了部署管道。本示例中，我们将演示一种设置超时的方法，可以针对每个测试设置不同的超时。

\mySmallsection{准备工作}

这个示例是一个Python脚本(\texttt{test.py})，它总是返回0。为了保持这种简单性，并保持对CMake方面的关注，测试脚本除了等待两秒钟外什么也不做。实际中，这个测试脚本将执行更有意义的工作:

\begin{codebox}[python]
import sys
import time

# wait for 2 seconds
time.sleep(2)

# report success
sys.exit(0)
\end{codebox}

\mySmallsection{具体实施}

我们需要通知CTest终止测试，如下:

\begin{enumerate}
\item 我们定义项目名称，启用测试，并定义测试:

\begin{codebox}[cmake]
# set minimum cmake version
cmake_minimum_required(VERSION 3.5 FATAL_ERROR)
# project name
project(recipe-07 LANGUAGES NONE)
# detect python
find_package(PythonInterp REQUIRED)
# define tests
enable_testing()
# we expect this test to run for 2 seconds
add_test(example ${PYTHON_EXECUTABLE} ${CMAKE_CURRENT_SOURCE_DIR}/test.py)
\end{codebox}

\item 另外，我们为测试指定时限，设置为10秒:

\begin{codebox}[cmake]
set_tests_properties(example PROPERTIES TIMEOUT 10)
\end{codebox}

\item 知道了如何进行配置和构建，并希望测试能够通过:

\begin{codebox}[shell]
$ ctest
Test project /home/user/cmake-recipes/chapter-04/recipe-07/example/build
Start 1: example
1/1 Test #1: example .......................... Passed 2.01 sec
100% tests passed, 0 tests failed out of 1
Total Test time (real) = 2.01 sec
\end{codebox}

\item 现在，为了验证超时是否有效，我们将\texttt{test.py}中的\texttt{sleep}命令增加到11秒，并重新运行测试:

\begin{codebox}[shell]
$ ctest
Test project /home/user/cmake-recipes/chapter-04/recipe-07/example/build
Start 1: example
1/1 Test #1: example ..........................***Timeout 10.01 sec
0% tests passed, 1 tests failed out of 1
Total Test time (real) = 10.01 sec
The following tests FAILED:
1 - example (Timeout)
Errors while running CTest
\end{codebox}

\end{enumerate}

\mySmallsection{工作原理}

\texttt{TIMEOUT}是一个方便的属性，可以使用\texttt{set\_tests\_properties}为单个测试指定超时时间。如果测试运行超过了这个设置时间，不管出于什么原因(测试已经停止或者机器太慢)，测试将被终止并标记为失败。


% --- recipe 8: 4.8-chinese.md ---
\section{4.8 并行测试}

\begin{myInfoBox}
\textbf{NOTE}：此示例代码可以在 \texthref{https://github.com/dev-cafe/cmake-cookbook/tree/v1.0/chapter-04/recipe-08} 中找到。该示例在CMake 3.5版(或更高版本)中是有效的，并且已经在GNU/Linux、macOS和Windows上进行过测试。
\end{myInfoBox}

大多数现代计算机都有4个或更多个CPU核芯。CTest有个非常棒的特性，能够并行运行测试，如果您有多个可用的核。这可以减少测试的总时间，而减少总测试时间才是真正重要的，从而开发人员频繁地进行测试。本示例中，我们将演示这个特性，并讨论如何优化测试以获得最大的性能。

其他测试可以进行相应地表示，我们把这些测试脚本放在\texttt{CMakeLists.txt}同目录下面的test目录中。

\mySmallsection{准备工作}

我们假设测试集包含标记为a, b，…，j的测试用例，每一个都有特定的持续时间:

\begin{table}[htbp]\centering
\begin{tabularx}{\textwidth}{lX}
\toprule
测试用例 & 该单元的耗时 \\
\midrule
a, b, c, d & 0.5 \\
e, f, g & 1.5 \\
h & 2.5 \\
i & 3.5 \\
j & 4.5 \\
\bottomrule
\end{tabularx}
\end{table}

时间单位可以是分钟，但是为了保持简单和简短，我们将使用秒。为简单起见，我们可以用Python脚本表示\texttt{test a}，它消耗0.5个时间单位:

\begin{codebox}[python]
import sys
import time

# wait for 0.5 seconds
time.sleep(0.5)

# finally report success
sys.exit(0)
\end{codebox}

其他测试同理。我们将把这些脚本放在\texttt{CMakeLists.txt}下面，一个名为\texttt{test}的目录中。

\mySmallsection{具体实施}

对于这个示例，我们需要声明一个测试列表，如下:

\begin{enumerate}
\item \texttt{CMakeLists.txt}非常简单：

\begin{codebox}[cmake]
# set minimum cmake version
cmake_minimum_required(VERSION 3.5 FATAL_ERROR)
# project name
project(recipe-08 LANGUAGES NONE)
# detect python
find_package(PythonInterp REQUIRED)
# define tests
enable_testing()
add_test(a ${PYTHON_EXECUTABLE} ${CMAKE_CURRENT_SOURCE_DIR}/test/a.py)
add_test(b ${PYTHON_EXECUTABLE} ${CMAKE_CURRENT_SOURCE_DIR}/test/b.py)
add_test(c ${PYTHON_EXECUTABLE} ${CMAKE_CURRENT_SOURCE_DIR}/test/c.py)
add_test(d ${PYTHON_EXECUTABLE} ${CMAKE_CURRENT_SOURCE_DIR}/test/d.py)
add_test(e ${PYTHON_EXECUTABLE} ${CMAKE_CURRENT_SOURCE_DIR}/test/e.py)
add_test(f ${PYTHON_EXECUTABLE} ${CMAKE_CURRENT_SOURCE_DIR}/test/f.py)
add_test(g ${PYTHON_EXECUTABLE} ${CMAKE_CURRENT_SOURCE_DIR}/test/g.py)
add_test(h ${PYTHON_EXECUTABLE} ${CMAKE_CURRENT_SOURCE_DIR}/test/h.py)
add_test(i ${PYTHON_EXECUTABLE} ${CMAKE_CURRENT_SOURCE_DIR}/test/i.py)
add_test(j ${PYTHON_EXECUTABLE} ${CMAKE_CURRENT_SOURCE_DIR}/test/j.py)
\end{codebox}

\item 我们可以配置项目，使用\texttt{ctest}运行测试，总共需要17秒:

\begin{codebox}[shell]
$ mkdir -p build
$ cd build
$ cmake ..
$ ctest
Start 1: a
1/10 Test #1: a ................................ Passed 0.51 sec
Start 2: b
2/10 Test #2: b ................................ Passed 0.51 sec
Start 3: c
3/10 Test #3: c ................................ Passed 0.51 sec
Start 4: d
4/10 Test #4: d ................................ Passed 0.51 sec
Start 5: e
5/10 Test #5: e ................................ Passed 1.51 sec
Start 6: f
6/10 Test #6: f ................................ Passed 1.51 sec
Start 7: g
7/10 Test #7: g ................................ Passed 1.51 sec
Start 8: h
8/10 Test #8: h ................................ Passed 2.51 sec
Start 9: i
9/10 Test #9: i ................................ Passed 3.51 sec
Start 10: j
10/10 Test #10: j ................................ Passed 4.51 sec
100% tests passed, 0 tests failed out of 10
Total Test time (real) = 17.11 sec
\end{codebox}

\item 现在，如果机器有4个内核可用，我们可以在不到5秒的时间内在4个内核上运行测试集:

\begin{codebox}[shell]
$ ctest --parallel 4
Start 10: j
Start 9: i
Start 8: h
Start 5: e
1/10 Test #5: e ................................ Passed 1.51 sec
Start 7: g
2/10 Test #8: h ................................ Passed 2.51 sec
Start 6: f
3/10 Test #7: g ................................ Passed 1.51 sec
Start 3: c
4/10 Test #9: i ................................ Passed 3.63 sec
5/10 Test #3: c ................................ Passed 0.60 sec
Start 2: b
Start 4: d
6/10 Test #6: f ................................ Passed 1.51 sec
7/10 Test #4: d ................................ Passed 0.59 sec
8/10 Test #2: b ................................ Passed 0.59 sec
Start 1: a
9/10 Test #10: j ................................ Passed 4.51 sec
10/10 Test #1: a ................................ Passed 0.51 sec
100% tests passed, 0 tests failed out of 10
Total Test time (real) = 4.74 sec
\end{codebox}

\end{enumerate}

\mySmallsection{工作原理}

可以观察到，在并行情况下，测试j、i、h和e同时开始。当并行运行时，总测试时间会有显著的减少。观察\texttt{ ctest --parallel 4}的输出，我们可以看到并行测试运行从最长的测试开始，最后运行最短的测试。从最长的测试开始是一个非常好的策略。这就像打包移动的盒子：从较大的项目开始，然后用较小的项目填补空白。a-j测试在4个核上的叠加比较，从最长的开始，如下图所示:

\begin{codebox}[shell]
--> time
core 1: jjjjjjjjj
core 2: iiiiiiibd
core 3: hhhhhggg
core 4: eeefffac
\end{codebox}

按照定义测试的顺序运行，运行结果如下:

\begin{codebox}[shell]
--> time
core 1: aeeeiiiiiii
core 2: bfffjjjjjjjjj
core 3: cggg
core 4: dhhhhh
\end{codebox}

按照定义测试的顺序运行测试，总的来说需要更多的时间，因为这会让2个核大部分时间处于空闲状态(这里的核3和核4)。CMake知道每个测试的时间成本，是因为我们先顺序运行了测试，将每个测试的成本数据记录在\texttt{test/Temporary/CTestCostData.txt}文件中:

\begin{codebox}[shell]
a 1 0.506776
b 1 0.507882
c 1 0.508175
d 1 0.504618
e 1 1.51006
f 1 1.50975
g 1 1.50648
h 1 2.51032
i 1 3.50475
j 1 4.51111
\end{codebox}

如果在配置项目之后立即开始并行测试，它将按照定义测试的顺序运行测试，在4个核上的总测试时间明显会更长。这意味着什么呢？这意味着，我们应该减少的时间成本来安排测试？这是一种决策，但事实证明还有另一种方法，我们可以自己表示每次测试的时间成本:

\begin{codebox}[cmake]
add_test(a ${PYTHON_EXECUTABLE} ${CMAKE_CURRENT_SOURCE_DIR}/test/a.py)
add_test(b ${PYTHON_EXECUTABLE} ${CMAKE_CURRENT_SOURCE_DIR}/test/b.py)
add_test(c ${PYTHON_EXECUTABLE} ${CMAKE_CURRENT_SOURCE_DIR}/test/c.py)
add_test(d ${PYTHON_EXECUTABLE} ${CMAKE_CURRENT_SOURCE_DIR}/test/d.py)
set_tests_properties(a b c d PROPERTIES COST 0.5)

add_test(e ${PYTHON_EXECUTABLE} ${CMAKE_CURRENT_SOURCE_DIR}/test/e.py)
add_test(f ${PYTHON_EXECUTABLE} ${CMAKE_CURRENT_SOURCE_DIR}/test/f.py)
add_test(g ${PYTHON_EXECUTABLE} ${CMAKE_CURRENT_SOURCE_DIR}/test/g.py)
set_tests_properties(e f g PROPERTIES COST 1.5)

add_test(h ${PYTHON_EXECUTABLE} ${CMAKE_CURRENT_SOURCE_DIR}/test/h.py)
set_tests_properties(h PROPERTIES COST 2.5)

add_test(i ${PYTHON_EXECUTABLE} ${CMAKE_CURRENT_SOURCE_DIR}/test/i.py)
set_tests_properties(i PROPERTIES COST 3.5)

add_test(j ${PYTHON_EXECUTABLE} ${CMAKE_CURRENT_SOURCE_DIR}/test/j.py)
set_tests_properties(j PROPERTIES COST 4.5)
\end{codebox}

成本参数可以是一个估计值，也可以从\texttt{test/Temporary/CTestCostData.txt}中提取。

\mySmallsection{更多信息}

除了使用\texttt{ ctest --parallel N}，还可以使用环境变量\texttt{CTEST\_PARALLEL\_LEVEL}将其设置为所需的级别。


% --- recipe 9: 4.9-chinese.md ---
\section{4.9 运行测试子集}

\begin{myInfoBox}
\textbf{NOTE}：此示例代码可以在 \texthref{https://github.com/dev-cafe/cmake-cookbook/tree/v1.0/chapter-04/recipe-09} 中找到。该示例在CMake 3.5版(或更高版本)中是有效的，并且已经在GNU/Linux、macOS和Windows上进行过测试。
\end{myInfoBox}

前面的示例中，我们学习了如何在CMake的帮助下并行运行测试，并讨论了从最长的测试开始是最高效的。虽然，这种策略将总测试时间最小化，但是在特定特性的代码开发期间，或者在调试期间，我们可能不希望运行整个测试集。对于调试和代码开发，我们只需要能够运行选定的测试子集。在本示例中，我们将实现这一策略。

\mySmallsection{准备工作}

在这个例子中，我们假设总共有六个测试：前三个测试比较短，名称分别为\texttt{feature-a}、\texttt{feature-b}和\texttt{feature-c}，还有三个长测试，名称分别是\texttt{feature-d}、\texttt{benchmark-a}和\texttt{benchmark-b}。这个示例中，我们可以用Python脚本表示这些测试，可以在其中调整休眠时间:

\begin{codebox}[python]
import sys
import time

# wait for 0.1 seconds
time.sleep(0.1)

# finally report success
sys.exit(0)
\end{codebox}

\mySmallsection{具体实施}

以下是我们CMakeLists.txt文件内容的详细内容:

\begin{enumerate}
\item \texttt{CMakeLists.txt}中，定义了六个测试:

\begin{codebox}[cmake]
cmake_minimum_required(VERSION 3.5 FATAL_ERROR)
# project name
project(recipe-09 LANGUAGES NONE)
# detect python
find_package(PythonInterp REQUIRED)
# define tests
enable_testing()
add_test(
  NAME feature-a
  COMMAND ${PYTHON_EXECUTABLE} ${CMAKE_CURRENT_SOURCE_DIR}/test/feature-a.py
  )
add_test(
  NAME feature-b
  COMMAND ${PYTHON_EXECUTABLE} ${CMAKE_CURRENT_SOURCE_DIR}/test/feature-b.py
  )
add_test(
  NAME feature-c
  COMMAND ${PYTHON_EXECUTABLE} ${CMAKE_CURRENT_SOURCE_DIR}/test/feature-c.py
  )
add_test(
  NAME feature-d
  COMMAND ${PYTHON_EXECUTABLE} ${CMAKE_CURRENT_SOURCE_DIR}/test/feature-d.py
  )
add_test(
  NAME benchmark-a
  COMMAND ${PYTHON_EXECUTABLE} ${CMAKE_CURRENT_SOURCE_DIR}/test/benchmark-a.py
  )
add_test(
  NAME benchmark-b
  COMMAND ${PYTHON_EXECUTABLE} ${CMAKE_CURRENT_SOURCE_DIR}/test/benchmark-b.py
  )
\end{codebox}

\item 此外，我们给较短的测试贴上\texttt{quick}的标签，给较长的测试贴上\texttt{long}的标签:

\begin{codebox}[cmake]
set_tests_properties(
  feature-a
  feature-b
  feature-c
  PROPERTIES
    LABELS "quick"
  )
set_tests_properties(
  feature-d
  benchmark-a
  benchmark-b
  PROPERTIES
    LABELS "long"
  )
\end{codebox}

\item 我们现在可以运行测试集了，如下:

\begin{codebox}[shell]
$ mkdir -p build
$ cd build
$ cmake ..
$ ctest
Start 1: feature-a
1/6 Test #1: feature-a ........................ Passed 0.11 sec
Start 2: feature-b
2/6 Test #2: feature-b ........................ Passed 0.11 sec
Start 3: feature-c
3/6 Test #3: feature-c ........................ Passed 0.11 sec
Start 4: feature-d
4/6 Test #4: feature-d ........................ Passed 0.51 sec
Start 5: benchmark-a
5/6 Test #5: benchmark-a ...................... Passed 0.51 sec
Start 6: benchmark-b
6/6 Test #6: benchmark-b ...................... Passed 0.51 sec
100% tests passed, 0 tests failed out of 6
Label Time Summary:
long = 1.54 sec*proc (3 tests)
quick = 0.33 sec*proc (3 tests)
Total Test time (real) = 1.87 sec
\end{codebox}

\end{enumerate}

\mySmallsection{工作原理}

现在每个测试都有一个名称和一个标签。CMake中所有的测试都是有编号的，所以它们也带有唯一编号。定义了测试标签之后，我们现在可以运行整个集合，或者根据它们的名称(使用正则表达式)、标签或编号运行测试。

按名称运行测试(运行所有具有名称匹配功能的测试):

\begin{codebox}[shell]
$ ctest -R feature

Start 1: feature-a
1/4 Test #1: feature-a ........................ Passed 0.11 sec
Start 2: feature-b
2/4 Test #2: feature-b ........................ Passed 0.11 sec
Start 3: feature-c
3/4 Test #3: feature-c ........................ Passed 0.11 sec
Start 4: feature-d
4/4 Test #4: feature-d ........................ Passed 0.51 sec
100% tests passed, 0 tests failed out of 4
\end{codebox}

按照标签运行测试(运行所有的长测试):

\begin{codebox}[shell]
$ ctest -L long

Start 4: feature-d
1/3 Test #4: feature-d ........................ Passed 0.51 sec
Start 5: benchmark-a
2/3 Test #5: benchmark-a ...................... Passed 0.51 sec
Start 6: benchmark-b
3/3 Test #6: benchmark-b ...................... Passed 0.51 sec
100% tests passed, 0 tests failed out of 3
\end{codebox}

根据数量运行测试(运行测试2到4)产生的结果是:

\begin{codebox}[shell]
$ ctest -I 2,4

Start 2: feature-b
1/3 Test #2: feature-b ........................ Passed 0.11 sec
Start 3: feature-c
2/3 Test #3: feature-c ........................ Passed 0.11 sec
Start 4: feature-d
3/3 Test #4: feature-d ........................ Passed 0.51 sec
100% tests passed, 0 tests failed out of 3
\end{codebox}

\mySmallsection{更多信息}

尝试使用\texttt{ \$ ctest --help }，将看到有大量的选项可供用来定制测试。


% --- recipe 10: 4.10-chinese.md ---
\section{4.10 使用测试固件}

\begin{myInfoBox}
\textbf{NOTE}：此示例代码可以在 \texthref{https://github.com/dev-cafe/cmake-cookbook/tree/v1.0/chapter-04/recipe-10} 中找到。该示例在CMake 3.5版(或更高版本)中是有效的，并且已经在GNU/Linux、macOS和Windows上进行过测试。
\end{myInfoBox}

这个示例的灵感来自于Craig Scott，我们建议读者也参考相应的博客文章来了解更多的背景知识，\texthref{https://crascit.com/2016/10/18/test-fixtures-withcmake-ctest/} ，此示例的动机是演示如何使用测试固件。这对于更复杂的测试非常有用，这些测试需要在测试运行前进行设置，以及在测试完成后执行清理操作(例如：创建示例数据库、设置连接、断开连接、清理测试数据库等等)。我们需要运行一个设置或清理操作的测试，并能够以一种可预测和健壮的方式自动触发这些步骤，而不需要引入代码重复。这些设置和清理步骤可以委托给测试框架(例如Google Test或Catch2)，我们在这里将演示如何在CMake级别实现测试固件。

\mySmallsection{准备工作}

我们将准备4个Python脚本，并将它们放在\texttt{test}目录下:\texttt{setup.py}、\texttt{features-a.py}、\texttt{features-b.py}和\texttt{clean-up.py}。

\mySmallsection{具体实施}

我们从\texttt{CMakeLists.txt}结构开始，附加一些步骤如下:

\begin{enumerate}
\item 基础CMake语句:

\begin{codebox}[cmake]
# set minimum cmake version
cmake_minimum_required(VERSION 3.5 FATAL_ERROR)
# project name
project(recipe-10 LANGUAGES NONE)
# detect python
find_package(PythonInterp REQUIRED)
# define tests
enable_testing()
\end{codebox}

\item 然后，定义了4个测试步骤，并将它们绑定到一个固件上:

\begin{codebox}[cmake]
add_test(
  NAME setup
  COMMAND ${PYTHON_EXECUTABLE} ${CMAKE_CURRENT_SOURCE_DIR}/test/setup.py
  )
set_tests_properties(
  setup
  PROPERTIES
    FIXTURES_SETUP my-fixture
  )
add_test(
  NAME feature-a
  COMMAND ${PYTHON_EXECUTABLE} ${CMAKE_CURRENT_SOURCE_DIR}/test/feature-a.py
  )
add_test(
  NAME feature-b
  COMMAND ${PYTHON_EXECUTABLE} ${CMAKE_CURRENT_SOURCE_DIR}/test/feature-b.py
  )
set_tests_properties(
  feature-a
  feature-b
  PROPERTIES
    FIXTURES_REQUIRED my-fixture
  )
add_test(
  NAME cleanup
  COMMAND ${PYTHON_EXECUTABLE} ${CMAKE_CURRENT_SOURCE_DIR}/test/cleanup.py
  )
set_tests_properties(
  cleanup
  PROPERTIES
    FIXTURES_CLEANUP my-fixture
  )
\end{codebox}

\item 运行整个集合，如下面的输出所示:

\begin{codebox}[shell]
$ mkdir -p build
$ cd build
$ cmake ..
$ ctest
Start 1: setup
1/4 Test #1: setup ............................ Passed 0.01 sec
Start 2: feature-a
2/4 Test #2: feature-a ........................ Passed 0.01 sec
Start 3: feature-b
3/4 Test #3: feature-b ........................ Passed 0.00 sec
Start 4: cleanup
4/4 Test #4: cleanup .......................... Passed 0.01 sec
100% tests passed, 0 tests failed out of 4
\end{codebox}

\item 然而，当我们试图单独运行测试特性时。它正确地调用设置步骤和清理步骤:

\begin{codebox}[shell]
$ ctest -R feature-a
Start 1: setup
1/3 Test #1: setup ............................ Passed 0.01 sec
Start 2: feature-a
2/3 Test #2: feature-a ........................ Passed 0.00 sec
Start 4: cleanup
3/3 Test #4: cleanup .......................... Passed 0.01 sec
100% tests passed, 0 tests failed out of 3
\end{codebox}

\end{enumerate}

\mySmallsection{工作原理}

在本例中，我们定义了一个文本固件，并将其称为\texttt{my-fixture}。我们为安装测试提供了\texttt{FIXTURES\_SETUP}属性，并为清理测试了\texttt{FIXTURES\_CLEANUP}属性，并且使用\texttt{FIXTURES\_REQUIRED}，我们确保测试\texttt{feature-a}和\texttt{feature-b}都需要安装和清理步骤才能运行。将它们绑定在一起，可以确保在定义良好的状态下，进入和离开相应的步骤。


